\section{Detalles de entrenamiento, costos y arquitecturas alternativas: por qué esto es ingeniería de sistemas}

Aunque la sección «Detalles de entrenamiento y costos» de la entrevista es breve, es la clave para entender todo el episodio. Luo Fuli (罗福莉) insiste en que, en la era de los Agent, todos los caminos deben volver a la tasa de finalización de extremo a extremo, la velocidad y el costo. Que un modelo sea más fuerte en un benchmark no significa que funcione mejor dentro de un framework de Agent; que un modelo sea más grande tampoco significa que valga la pena invocarlo en todas las etapas. La verdadera ingeniería de sistemas consiste en reunir en una misma tabla la capacidad del modelo, las etapas de la tarea, las llamadas a herramientas, la latencia, el precio y la confiabilidad para tomar decisiones.

\begin{importantbox}{El costo no es una partida financiera secundaria, sino parte de la capacidad del modelo}
Si una tarea de Agent requiere decenas de rondas de llamadas al modelo, ejecución de herramientas y retroalimentación del entorno, el costo de cada llamada se multiplica a lo largo de la tarea de horizonte largo. Los modelos pequeños y baratos, los modelos especializados, el enrutamiento entre varios modelos, la caché, los verificadores y la orquestación del framework influyen directamente en que el sistema pueda ponerse en producción.
\end{importantbox}

\subsection{Términos clave: vocabulario de entrenamiento y costos}

\noindent\begin{tabularx}{\textwidth}{p{0.23\textwidth}p{0.31\textwidth}X}
\toprule
Término & Problema que resuelve & Relación con este episodio \\
\midrule
Tasa de finalización de extremo a extremo & Si la tarea completa termina cumpliéndose & Un producto de Agent no puede evaluarse solo por acciones parciales o respuestas de una sola ronda. \\
Enrutamiento de llamadas & Elegir modelos distintos para etapas distintas & Usar un modelo grande para el razonamiento crítico y uno pequeño para las etapas frecuentes y de bajo riesgo. \\
Verificador & Determinar si el resultado de la tarea es correcto & Es esencial para el RL feedback y la evaluación de confiabilidad. \\
Costo de Rollout & Costo de cómputo y de API de generar múltiples trayectorias & Determina si el postentrenamiento puede hacer scale. \\
Latency & Tiempo de espera del usuario y duración total de la tarea & Influye en que una herramienta de productividad sea utilizable. \\
Scaffold & Framework externo del Agent, orquestación de herramientas y gestión de estado & Puede compensar las deficiencias del modelo y también elevar el techo de los modelos de punta. \\
\bottomrule
\end{tabularx}

\subsection{Por qué importan las arquitecturas alternativas}

Aquí las arquitecturas alternativas no son «trucos ingeniosos sin sustancia», sino investigación escalable bajo restricciones de recursos. DeepSeek-V2, MoE, las modificaciones de attention y otras líneas de trabajo plantean, en esencia, la misma pregunta: cómo obtener más inteligencia por unidad de costo con un cómputo y un hardware dados. En la era de los Agent esta pregunta se vuelve más apremiante, porque las tareas de horizonte largo invocan el modelo una y otra vez, y cualquier diferencia de eficiencia en un solo paso se amplifica.

\begin{warningbox}{No hay que ver las arquitecturas alternativas solo como innovación de artículos}
Si una arquitectura no puede hacer scale hasta un modelo industrial, no puede desplegarse en un sistema de inferencia real o no logra reducir el costo de extremo a extremo de las tareas, su valor para los productos de Agent es limitado. El «nivel de calidad industrial» que subraya la entrevista exige justamente que la innovación arquitectónica termine incorporándose a sistemas que funcionen.
\end{warningbox}

\subsection{Resumen de la sección}

Los detalles de entrenamiento, los costos y las arquitecturas alternativas muestran en conjunto que la ruta tecnológica de la era de los Agent no es una competencia de una sola variable por «modelos más grandes», sino una competencia por la eficiencia del sistema de extremo a extremo. La estructura del modelo, el postentrenamiento, el framework y el costo de inferencia deben optimizarse en conjunto.
